%% ======================================================================
\section{From the Gibbs classifier to the majority vote}
\label{chap:mv}
%% ======================================================================

\begin{objectives}
\guiding{Why is the vote better than its voters, and how can we certify the vote
itself rather than a voter drawn at random?}
The bounds of \cref{chap:bounds} concern the Gibbs classifier, whose risk is the
risk of an average voter. Ensemble methods, however, predict with the majority
vote, whose whole point is to be better than the average voter. This section
relates the two. We show that the risk of the vote is a tail probability of the
proportion of erring voters, and we control this tail with moments, which leads
to the first-order bound, the tandem bound and the C-bound. We compare them with
Condorcet's jury theorem and on random forests, prove which one is the tightest
in which regime, and show how PAC-Bayes estimates each of them from data.
\end{objectives}

Throughout this section, the voters take values in $\{-1,+1\}$ and we use the
0-1 loss.

\subsection{Two random variables that describe the vote}

To compare the vote with its voters, we need a common language. A single
random variable will do: on each example, the proportion of voters that err.
For an example $(x,y)$, let
\begin{align*}
W_Q(x,y)&=\Esp{h\sim Q}{\ind{h(x)\neq y}}\in[0,1],\\
M_Q(x,y)&=y\,\Esp{h\sim Q}{h(x)}=1-2W_Q(x,y)\in[-1,1].
\end{align*}
The variable $W_Q$ is the $Q$-weight of the voters that make a mistake on
$(x,y)$, and $M_Q$ is the \emph{margin} of the vote, already met in
\cref{chap:intro} for AdaBoost. The vote errs exactly when at least half of the
weight is on the wrong side, so that, with our convention for ties,
\begin{equation}\label{eq:mv_risk_W}
\RD(\BQ)=\P_{(x,y)\sim\D}\big(M_Q(x,y)\leq0\big)=\P_{(x,y)\sim\D}\Big(W_Q(x,y)\geq\tfrac12\Big),
\end{equation}
whereas the Gibbs risk is its expectation, $\RD(\GQ)=\E_{\D}W_Q$. In other words,
the risk of the majority vote is a \emph{tail probability} of the random
variable $W_Q$, and the Gibbs risk is its \emph{first moment}. All the bounds of
this section control this tail with moments of $W_Q$, exactly as Markov's,
Chebyshev's and Cantelli's inequalities control the tail of a random variable
with its first two moments.

\Cref{fig:wq} shows the distribution of $W_Q$ for a random forest of $200$ trees
on the \texttt{digits} dataset (digits $0$--$4$ against $5$--$9$). This forest is
larger than the one of our running example and is evaluated on a different
split, hence the slightly different values. Most test points
are classified correctly by a large majority of trees, so that the histogram is
concentrated near zero; its mean, the Gibbs risk, is $0.164$. Only the points in
the red area, where at least half of the trees err, are misclassified by the
vote, and they represent $3.5\%$ of the test sample. The vote is thus almost
five times more accurate than an average tree, and a good bound on the vote must
be able to see this. Let us check how far each moment inequality can go,
starting with the crudest one.

\begin{figure}[t]
\centering
\includegraphics[width=\linewidth]{wq_distribution}
\caption{Random forest of 200 trees on \texttt{digits} (classes $0$--$4$ vs
$5$--$9$). Left: histogram of $W_Q$ on the test set; the Gibbs risk is the mean
of this distribution ($0.164$), while the majority vote only errs when
$W_Q\geq1/2$ (red area, $3.5\%$ of the points). Right: distribution of the
margins $M_Q=1-2W_Q$.}
\label{fig:wq}
\end{figure}

\subsection{The first-order bound}

Since we already know how to bound the Gibbs risk, the first moment of $W_Q$,
the simplest idea is to turn it into a bound on the tail with Markov's
inequality.

\begin{proposition}{First-order bound (folklore, e.g.\ \citealp{langford2002pac})}{fo}
For any $Q\in\M(\Hcal)$, $\RD(\BQ)\leq2\,\RD(\GQ)$.
\end{proposition}
\begin{proof}
By \eqref{eq:mv_risk_W} and Markov's inequality,
$\RD(\BQ)=\P(W_Q\geq\frac12)\leq\frac{\E W_Q}{1/2}=2\RD(\GQ)$.
\end{proof}

In words, the vote is never worse than twice an average voter: whatever the
ensemble, voting cannot go badly wrong. Combined with Seeger's bound, the
first-order bound gives a first PAC-Bayesian
certificate for the majority vote: with probability at least $1-\delta$ over
$S\sim\D^m$, for all $Q$,
\begin{equation}\label{eq:fo_pb}
\RD(\BQ)\leq2\,\klup\Big(\RS(\GQ),\frac{\KL(Q\|P)+\ln\frac{2\sqrt m}{\delta}}{m}\Big).
\end{equation}
The factor $2$ cannot be improved in general: if $W_Q$ only takes the values
$0$ and $\frac12$, the vote errs on every example where $W_Q=\frac12$ and the
inequality is an equality. But this worst case is precisely the one where voting
is useless, and the first-order bound completely ignores the benefit of the
vote. On the forest of \cref{fig:wq}, it gives $0.327$, twice the risk of an
average tree, for a vote whose risk is ten times smaller ($0.035$). Worse, as
we will see in \cref{sec:algo_fo}, \emph{minimizing} this bound pushes $Q$
towards the few best voters and destroys the diversity of the ensemble. To do
better, we must understand where the benefit of voting comes from.

\subsection{Condorcet's jury theorem}

To understand what a good bound on the vote should capture, consider the
most favourable situation, known since Condorcet (1785): the voters err
\emph{independently}.

\begin{example}[Condorcet's jury theorem]
Let $Q$ be uniform over $n=2k+1$ voters that err independently, each with
probability $p<1/2$. Then $nW_Q\sim\mathrm{Bin}(n,p)$, the Gibbs risk is $p$ and
$\RD(\BQ)=\P(\mathrm{Bin}(n,p)\geq k+1)$, which tends to $0$ exponentially fast
when $n\to\infty$. With $n=3$ and $p=0.3$, $\RD(\BQ)=3p^2(1-p)+p^3=0.216$; with
$n=25$, it is already $0.017$, and with $n=101$ it is about $10^{-5}$.
\end{example}

The left panel of \cref{fig:condorcet} plots this risk together with the bounds
of this section. The first-order bound stays at $2p=0.6$ whatever the number of
voters. The tandem bound, introduced below, improves on it but converges to
$4p^2=0.36$. Only the C-bound follows the risk of the vote down to zero, at the
rate $1/n$.

Independence is of course unrealistic. In the right panel, the $25$
voters copy the errors of a common ``leader'' with probability $c$ and err
independently otherwise, so that $c$ measures the correlation of their errors.
When $c$ increases, the vote loses its advantage, its risk goes from $0.018$ to
$0.30$, and all the bounds deteriorate; for strongly correlated voters, the
first-order bound becomes the best of the three. The rest of this section
explains these observations.

\begin{figure}[t]
\centering
\includegraphics[width=\linewidth]{condorcet}
\caption{Risk of the majority vote and oracle bounds for voters of individual
error $p=0.3$. Left: independent errors, as a function of the number of voters.
Right: $25$ voters whose errors are correlated (each voter copies the errors of
a common leader with probability $c$); $20\,000$ simulated examples.}
\label{fig:condorcet}
\end{figure}

\subsection{Second-order quantities: joint error and disagreement}

Condorcet's example contains the key idea: what makes a vote good is not only
the quality of each voter, but the fact that the voters do not all make their
mistakes on the same examples. A single voter cannot tell us this; we need to
look at \emph{pairs} of voters $(h,h')\sim Q^2=Q\otimes Q$ drawn independently
from the posterior.

\begin{definition}{Joint error and disagreement}{joint}
The \textbf{joint error} (or \emph{tandem loss}) and the \textbf{disagreement} of $Q$ are
\begin{align*}
\eD(Q)&=\Esp{(h,h')\sim Q^2}{\P_{(x,y)\sim\D}\big(h(x)\neq y\ \wedge\ h'(x)\neq y\big)}=\E_\D\big[W_Q^2\big],\\
\dD(Q)&=\Esp{(h,h')\sim Q^2}{\P_{x\sim\D_\X}\big(h(x)\neq h'(x)\big)}=\E_\D\big[2W_Q(1-W_Q)\big],
\end{align*}
and their empirical counterparts $\eS(Q)$, $\dS(Q)$ are defined on $S$.
\end{definition}

The expressions in terms of $W_Q$ follow from Fubini's theorem: for a fixed
example $(x,y)$, two independent draws $h,h'\sim Q$ both err with probability
$W_Q^2$, and exactly one of them errs with probability $2W_Q(1-W_Q)$, which,
for binary voters, is exactly the probability that they disagree. The joint
error is therefore the second moment of $W_Q$, and the Gibbs risk, the joint
error and the disagreement are linked by a simple identity.

\begin{proposition}{Decomposition of the Gibbs risk}{decomp}
For any $Q$,
\[
\RD(\GQ)=\eD(Q)+\tfrac12\dD(Q),\qquad\text{and}\qquad
0\leq\dD(Q)\leq2\RD(\GQ)\big(1-\RD(\GQ)\big).
\]
\end{proposition}
\begin{proof}
Pointwise, $W_Q=W_Q^2+\frac12\cdot2W_Q(1-W_Q)$; take expectations. For the
inequality, $\eD(Q)=\E W_Q^2\geq(\E W_Q)^2=\RD(\GQ)^2$ by Jensen's inequality, hence
$\dD(Q)=2(\RD(\GQ)-\eD(Q))\leq2\RD(\GQ)(1-\RD(\GQ))$.
\end{proof}

This decomposition says that, for a fixed Gibbs risk, a larger disagreement
means a smaller joint error: the voters make their mistakes on \emph{different}
examples, which is exactly what makes a vote effective. On the forest of
\cref{fig:wq}, the disagreement is $0.236$ and the joint error is only
$0.0457$. Note also that the disagreement does not depend on the labels: it can
be estimated on unlabeled data, which are often more plentiful, an observation
exploited in semi-supervised learning and domain adaptation
\citep{germain2020pac}. With these second-order quantities in hand, we can
revisit Markov's inequality.

\subsection{The tandem bound}

The first-order bound was loose because it only used the mean of $W_Q$.
Applying Markov's inequality to $W_Q^2$ instead of $W_Q$ gives a bound based on
the second moment, which sees the joint errors.

\begin{theorem}{Tandem (second-order) bound \citep{masegosa2020second}}{tnd}
For any $Q\in\M(\Hcal)$, $\RD(\BQ)\leq4\,\eD(Q)$.
\end{theorem}
\begin{proof}
By Markov's inequality applied to the non-negative variable $W_Q^2$,
$\P(W_Q\geq\frac12)=\P(W_Q^2\geq\frac14)\leq4\E W_Q^2=4\eD(Q)$.
\end{proof}

In words, wherever the vote errs, at least half of the weight is on the wrong
side, so two voters drawn independently from $Q$ both err there with probability
at least $\frac14$; the frequency of such joint errors therefore controls the
risk of the vote. On the forest of
\cref{fig:wq}, the tandem bound is $4\times0.0457=0.183$, almost twice smaller
than the first-order bound. By Proposition~\ref{prop:decomp},
$4\eD(Q)=2\RD(\GQ)-2\big(\dD(Q)-\RD(\GQ)\big)$, so the tandem bound improves on
the first-order bound exactly when $\dD(Q)\geq\RD(\GQ)$, that is when the
ensemble is diverse enough. In the worst case of identical voters, $\dD(Q)=0$ and
the tandem bound is twice the first-order bound.

To turn the tandem bound into a certificate, note that the joint error is the
Gibbs risk of the \emph{pair voter} $(h,h')\mapsto$ ``both err'' under the
posterior $Q^2$ on $\Hcal^2$. Applying Seeger's bound on $\Hcal^2$ with prior
$P^2$, and using $\KL(Q^2\|P^2)=2\KL(Q\|P)$, we obtain the following corollary.

\begin{corollary}{PAC-Bayesian tandem bound \citep{masegosa2020second}}{tnd_pb}
With probability at least $1-\delta$ over $S\sim\D^m$, for all $Q\in\M(\Hcal)$,
\begin{equation}\label{eq:tnd_pb}
\RD(\BQ)\leq4\,\klup\Big(\eS(Q),\ \frac{2\KL(Q\|P)+\ln\frac{2\sqrt m}{\delta}}{m}\Big).
\end{equation}
\end{corollary}

The price of the second-order information is a doubled KL term, and a factor
$4$ instead of $2$ in front of the kl inversion. The tandem certificate
therefore pays off when the ensemble is diverse and the sample is large enough
for the complexity term to be small.

\subsection{The C-bound}

The first-order bound uses the first moment of $W_Q$, the tandem bound its
second moment. Using both at once should do better, and this is what the
C-bound achieves through the Chebyshev--Cantelli inequality
(Lemma~\ref{lem:cantelli}).

\begin{theorem}{C-bound \citep{lacasse2007pac,germain2015risk}}{cbound}
For any $Q$ such that $\RD(\GQ)<\frac12$ (equivalently $\E_\D M_Q>0$),
\begin{equation}\label{eq:cbound}
\RD(\BQ)\leq\mathcal C_Q:=1-\frac{\big(\E_\D M_Q\big)^2}{\E_\D M_Q^2}
=1-\frac{\big(1-2\RD(\GQ)\big)^2}{1-2\dD(Q)}.
\end{equation}
\end{theorem}
\begin{proof}
Let $\mu=\E M_Q>0$. By Cantelli's inequality with $a=\mu$,
$\RD(\BQ)=\P(M_Q\leq0)=\P(M_Q-\mu\leq-\mu)\leq\frac{\Var M_Q}{\Var M_Q+\mu^2}
=\frac{\E M_Q^2-\mu^2}{\E M_Q^2}$.
Finally, $\E M_Q=1-2\E W_Q=1-2\RD(\GQ)$ and
$\E M_Q^2=\E(1-2W_Q)^2=1-4\E W_Q+4\E W_Q^2=1-2\dD(Q)$ by Proposition~\ref{prop:decomp}.
\end{proof}

The first expression of the C-bound has a nice interpretation: it is small when
the margin is large on average and has a small variance, i.e.\ when the vote is
correct with a comfortable and stable majority. The second expression, in terms
of the Gibbs risk and the disagreement, makes the classical intuition on ensembles
precise; \cref{fig:cbound} shows its landscape. For a fixed Gibbs risk, the
C-bound decreases when the disagreement increases: \emph{a good vote needs
voters that are individually better than random and that disagree}. This is
exactly what bagging and random forests try to achieve by randomizing the
training of the trees, and what \citet{breiman2001random} described with the
\emph{strength} and the \emph{correlation} of the trees. On the forest of
\cref{fig:wq}, the C-bound is $0.143$ (with the unrounded values of the Gibbs
risk and of the disagreement), the smallest of the three bounds. Is this always
the case? The next subsection answers precisely.

\begin{figure}[t]
\centering
\includegraphics[width=0.62\linewidth]{cbound_landscape}
\caption{Value of the C-bound as a function of the Gibbs risk and the
disagreement, on the feasible domain $\dD\leq2R(1-R)$. For a fixed Gibbs risk,
increasing the disagreement decreases the bound: the C-bound formalizes the
\emph{strength--diversity} trade-off of ensemble methods \citep{breiman2001random}.}
\label{fig:cbound}
\end{figure}

\subsection{Which bound is the tightest?}

The three oracle bounds only depend on the pair $(\RD(\GQ),\dD(Q))$, since
$4\eD(Q)=4\RD(\GQ)-2\dD(Q)$. They can therefore be compared once and for all,
and the comparison turns out to be remarkably simple.

\begin{proposition}{Comparison of the oracle bounds}{compare}
Let $R=\RD(\GQ)<\frac12$ and $d=\dD(Q)$. Then the C-bound is never larger than
the tandem bound, $\mathcal C_Q\leq4\eD(Q)$, and
\[
d\geq R\ \Longrightarrow\ \mathcal C_Q\leq4\eD(Q)\leq2R,\qquad\qquad
d\leq R\ \Longrightarrow\ 2R\leq\mathcal C_Q\leq4\eD(Q).
\]
The three bounds are equal on the line $d=R$.
\end{proposition}
\begin{proof}
The inequality $\mathcal C_Q\leq 4\eD(Q)$ follows from Theorem~\ref{thm:cctnd}
below: the tandem bound is the member $\mu=0$ of a family whose minimum is the
C-bound. We have seen that $4\eD(Q)\leq2R$ if and only if $d\geq R$. For the
C-bound, since $1-2R>0$ and $1-2d>0$,
\[
\mathcal C_Q\leq2R\iff\frac{(1-2R)^2}{1-2d}\geq1-2R\iff1-2R\geq1-2d\iff d\geq R.
\]
On the line $d=R$, all three expressions equal $2R$.
\end{proof}

\begin{figure}[t]
\centering
\includegraphics[width=0.58\linewidth]{bound_regions}
\caption{The smallest of the three oracle bounds in the plane (Gibbs risk,
disagreement). Above the line $d_Q=R(G_Q)$, the C-bound is the smallest; below,
the first-order bound is. The tandem bound is never strictly the smallest.}
\label{fig:bound_regions}
\end{figure}

In plain words: when two random voters disagree more often than an average
voter errs, the C-bound is the best of the three; otherwise, the simple
first-order bound wins. \Cref{fig:bound_regions} draws the two regions.

The proposition explains the observations made on Condorcet's example:
independent voters have a large
disagreement, $d=2p(1-p)(1-1/n)$, which exceeds $p$ as soon as $n\geq5$ when
$p=0.3$, so that the C-bound is the best; strongly correlated voters rarely
disagree, and the first-order bound becomes the best. The proposition also shows
that, as \emph{oracle} quantities, the C-bound dominates the tandem bound. The
situation changes once the bounds have to be \emph{estimated} from data: the tandem
bound needs a single kl inversion on one quantity, while the C-bound needs two,
and the ratio in its definition amplifies the estimation errors. In practice
the two certificates are often close, as we will see in \cref{chap:algos}.

\paragraph{PAC-Bayesian C-bound.} Let us see how the C-bound is estimated.
Since $\mathcal C_Q$ is increasing in
$\RD(\GQ)$ on $[0,\frac12)$ and decreasing in $\dD(Q)$, an upper bound on
the first and a lower bound on the second give an upper bound on $\mathcal C_Q$.
With a union bound (confidence $\delta/2$ for each), with probability at least
$1-\delta$, for all $Q$,
\begin{equation}\label{eq:cbound_pb}
\RD(\BQ)\leq1-\frac{\big(1-2\,\overline r\big)^2}{1-2\,\underline d},\qquad\text{with}
\end{equation}
\[
\overline r=\klup\Big(\RS(\GQ),\tfrac{\KL(Q\|P)+\ln\frac{4\sqrt m}{\delta}}{m}\Big),\qquad
\underline d=\kllow\Big(\dS(Q),\tfrac{2\KL(Q\|P)+\ln\frac{4\sqrt {m'}}{\delta}}{m'}\Big),
\]
provided $\overline r<\frac12$, where $m'$ is the number of (possibly unlabeled)
examples used to estimate the disagreement. \citet{lacasse2007pac} and
\citet{germain2015risk} give tighter versions that bound the pair
$(\RD(\GQ),\dD(Q))$ jointly, and \citet{viallard2021self} derive C-bounds that
can be minimized directly by gradient descent.

\subsection{A unifying view: the Chebyshev--Cantelli family}\label{sec:cctnd}

We have met three bounds, each obtained from $\P(W_Q\geq\frac12)$ with a
moment inequality. Are they really different objects? \citet{wu2021chebyshev}
noticed that the last two belong to a single parametric family.

\begin{theorem}{Parametric second-order bound \citep{wu2021chebyshev}}{cctnd}
For any $Q$ and any $\mu<\frac12$,
\begin{equation}\label{eq:cctnd}
\RD(\BQ)\leq\frac{\E_\D\big[(W_Q-\mu)^2\big]}{\big(\frac12-\mu\big)^2}
=\frac{\eD(Q)-2\mu\RD(\GQ)+\mu^2}{\big(\frac12-\mu\big)^2}.
\end{equation}
For $\mu=0$ it is the tandem bound, and, when $\RD(\GQ)<\frac12$, its minimum over
$\mu<\frac12$ is reached at $\mu^\star=\frac{\RD(\GQ)/2-\eD(Q)}{1/2-\RD(\GQ)}$ and
equals the C-bound $\mathcal C_Q$.
\end{theorem}
\begin{proof}
For $\mu<\frac12$, $\{W_Q\geq\frac12\}\subseteq\{(W_Q-\mu)^2\geq(\frac12-\mu)^2\}$;
apply Markov's inequality. The minimization over $\mu$ is left as an exercise
(TD, Exercise~8); it reproduces the classical proof of Cantelli's inequality.
\end{proof}

In words, instead of measuring how far $W_Q$ is from $0$, as the tandem bound
does, we measure how far it is from a well-chosen centre $\mu$; the best centre
gives back the C-bound. \Cref{fig:cctnd} draws this family on a random forest of $100$ trees trained on
a simulated problem, for which the Gibbs risk is $0.247$, the disagreement
$0.314$ and the joint error $0.090$. The curve starts from the tandem bound
$0.359$ at $\mu=0$, decreases to the C-bound $0.310$ at $\mu^\star=0.13$, and
blows up when $\mu$ approaches $\frac12$. The vote itself has a risk of $0.083$:
even the best member of the family is loose, which is the price of using only
two moments of $W_Q$.

\begin{figure}[t]
\centering
\includegraphics[width=0.66\linewidth]{cctnd}
\caption{The Chebyshev--Cantelli family \eqref{eq:cctnd} as a function of $\mu$,
for a random forest of $100$ trees on a simulated problem (oracle values
computed on the test sample).}
\label{fig:cctnd}
\end{figure}

If the minimum is the C-bound, why keep the whole family? The interest of the
parametric form is statistical. The quantity
$\E(W_Q-\mu)^2$ has a small variance when $\mu$ is well chosen, and it can be
estimated with a PAC-Bayes-Bennett inequality that exploits this small
variance. The resulting \emph{CCPBB} bound was the tightest of the bounds of this
section in the experiments of \citet{wu2021chebyshev} when enough data are
available. Related refinements exploit the fact that the excess losses involved
take only three values \citep[split-kl inequality,][]{wu2022split}.

\subsection{Margins and boosting}\label{sec:margins}

All the bounds so far summarize the margin distribution by one or two moments.
The \emph{margin theory} of boosting \citep{schapire1998boosting}, which we met
in \cref{fig:adaboost}, uses the whole distribution instead, and
Equation~\eqref{eq:mv_risk_W} is the bridge between the two viewpoints. For
any $\theta>0$, $\RD(\BQ)=\P(M_Q\leq0)\leq\P(M_Q\leq\theta)$, and the classical
margin bound for a finite class of voters states that, with probability at least
$1-\delta$,
\[
\P_\D(M_Q\leq0)\leq\P_S(M_Q\leq\theta)+O\left(\sqrt{\frac{\ln|\Hcal|\,\ln m}{m\,\theta^2}+\frac{\ln(1/\delta)}{m}}\right).
\]
This explains why AdaBoost may keep improving after its training error has
reached zero: it keeps increasing the margins, and the empirical term
$\P_S(M_Q\leq\theta)$ decreases for every $\theta$.

A PAC-Bayesian proof of such
bounds is obtained by approximating $\BQ$ by a vote of $N$ voters drawn from $Q$
\citep{langford2002pac,mcallester2003simplified}, and \citet{biggs2022margins}
sharpened margin bounds for voting classifiers using Dirichlet posteriors. The
C-bound can be read as a margin bound that uses the first two moments of the
margin distribution instead of its whole empirical distribution function.

\subsection{Back to the running example}

We now have all the tools to return to the two questions raised by our forest
in \cref{chap:intro}.

\begin{running}
We can now answer the first question of \cref{chap:intro}. On the test sample,
the trees of our forest have a Gibbs risk of $0.190$ and a disagreement of
$0.266$: two trees drawn at random disagree on more than a quarter of the digits,
and they both err on only $5.7\%$ of them. Since the disagreement exceeds the
Gibbs risk, Proposition~\ref{prop:compare} tells us that the C-bound is the
smallest oracle bound, and indeed the first-order bound is $0.381$, the tandem
bound $0.229$ and the C-bound $0.179$, for a vote whose risk is $0.044$. The
vote is good \emph{because the trees are diverse}: they are individually
mediocre, but their errors are spread over different digits.

The second question, certifying the vote without a test set, receives a more
modest answer. Estimating the three quantities on the out-of-bag examples (as in
the running example of \cref{chap:bounds}), with $\delta=0.05$, gives the
certificates $0.460$ (first order), $0.387$ (tandem) and $0.559$ (C-bound).
They are valid, but far from the test risk: the factor $4$ of the tandem bound
and the doubled KL term are expensive, and the C-bound suffers from its two kl
inversions. Improving these certificates by \emph{learning} the weights of the
trees is the program of Part~III.
\end{running}

\subsection{Summary and illustration on random forests}

To close the section, \cref{fig:mv_diversity} puts all the pieces together on
random forests of $100$ trees, where we vary the number of features tried at each split
(\texttt{max\_features}), the main diversity parameter of the algorithm. With a
single feature per split, the trees are weak (Gibbs risk $0.30$) but very diverse
(disagreement $0.38$); with all $20$ features, they are stronger (Gibbs risk
$0.17$) but much more similar (disagreement $0.19$). The risk of the vote is
smallest in between, around $4$ to $8$ features, where it reaches $0.077$.

The first-order bound only sees the strength of the trees and decreases monotonically,
from $0.60$ to $0.34$, missing the optimum. The tandem bound and the C-bound see
both strength and diversity: they are minimal at $8$ features, close to the
best vote, with values $0.29$ and $0.28$. Since the disagreement always exceeds
the Gibbs risk in this experiment, Proposition~\ref{prop:compare} predicts the
order C-bound $\leq$ tandem $\leq$ first order, which is what we observe.

\begin{figure}[t]
\centering
\includegraphics[width=\linewidth]{mv_bounds_diversity}
\caption{Random forests of 100 trees on a simulated problem, for an increasing
number of features tried at each split (\texttt{max\_features}). Left: oracle
values of the bounds computed on a large test sample. Right: stronger trees
(small Gibbs risk) are also less diverse (small disagreement). The C-bound
and the tandem bound track the risk of the vote much better than the
first-order bound.}
\label{fig:mv_diversity}
\end{figure}

The following table summarizes the bounds of the section: which moment of $W_Q$
they use, and how PAC-Bayes estimates them.

\begin{center}
\renewcommand{\arraystretch}{1.35}
\small
\begin{tabular}{llll}
\toprule
Bound & Oracle form & Moment used & PAC-Bayesian estimation \\
\midrule
First order & $2\RD(\GQ)$ & $\E W_Q$ & Seeger on $\RS(\GQ)$, $\KL$ \\
Tandem & $4\eD(Q)$ & $\E W_Q^2$ & Seeger on $\eS(Q)$, $2\KL$ \\
C-bound & $1-\frac{(1-2\RD(\GQ))^2}{1-2\dD(Q)}$ & $\E W_Q$, $\E W_Q^2$ & two kl inversions \\
Chebyshev--Cantelli & $\frac{\eD-2\mu\RD(\GQ)+\mu^2}{(1/2-\mu)^2}$ & $\E(W_Q-\mu)^2$ & PAC-Bayes-Bennett \\
\bottomrule
\end{tabular}
\end{center}

\begin{keypoints}
The risk of a vote is the tail probability $\P(W_Q\geq\frac12)$ of the
proportion of erring voters, whereas the Gibbs risk is its mean. The average
quality of the voters is not enough to control this tail: one also needs to
know how their errors are correlated, which is measured by the disagreement,
through the identity $\RD(\GQ)=\eD(Q)+\frac12\dD(Q)$. The first-order bound
$2\RD(\GQ)$ ignores the correlations, the tandem bound $4\eD(Q)$ and the C-bound
take them into account, and the C-bound, which is the minimum of the
Chebyshev--Cantelli family, is the smallest as soon as $\dD(Q)\geq\RD(\GQ)$.
PAC-Bayes estimates the second-order quantities on pairs of voters, with a
doubled KL term.
\end{keypoints}

\begin{quickchecks}
\item Express the risk of the vote and the Gibbs risk in terms of the random variable $W_Q$.
\item Why does the first-order bound ignore the benefit of voting?
\item Which quantity of the ensemble can be estimated without labels?
\item In which regime is the first-order bound smaller than the C-bound?
\end{quickchecks}

\begin{exercises}
\begin{exercise}
Give an example of three voters and a distribution on three points for which
$\RD(\BQ)=0$ while $\RD(\GQ)=\frac13$. Compute the three oracle bounds.
\end{exercise}
\begin{exercise}
Show that the C-bound is always smaller than $1$ when $\RD(\GQ)<\frac12$, and that it
equals $4\RD(\GQ)(1-\RD(\GQ))$ when $\dD(Q)=0$. Compare with the first-order bound.
\end{exercise}
\begin{exercise}
$(\star)$ Prove that the minimum over $\mu$ of the Chebyshev--Cantelli family equals the C-bound
(TD, Exercise~8).
\end{exercise}
\begin{exercise}
In Condorcet's model with $n$ independent voters of error $p$, compute $\eD(Q)$
and $\dD(Q)$, and show that the C-bound decreases like $1/n$ while the tandem
bound tends to $4p^2$.
\end{exercise}
\begin{exercise}
\textbf{(Comprehension)} True or false? Justify briefly. (a)~If $\RD(\GQ)<\frac12$, then
$\RD(\BQ)<\frac12$. (b)~If $\dD(Q)=0$, all the voters in the support of $Q$ make the same
predictions on almost every $x$. (c)~Estimating the joint error $\eD(Q)$ requires labels,
but estimating the disagreement $\dD(Q)$ does not. (d)~For a fixed Gibbs risk, the
tandem bound decreases when the disagreement increases.
\end{exercise}
\begin{exercise}
\textbf{(Application)} Three voters are evaluated on five equally likely examples
$z_1,\dots,z_5$. Voter $h_1$ errs only on $z_1$, $h_2$ errs on $z_2$ and $z_4$, and $h_3$ errs
on $z_2$ and $z_3$; the posterior is $Q=(0.4,0.35,0.25)$. Compute $W_Q(z_j)$ and $M_Q(z_j)$
for each example, the risk of the vote, the Gibbs risk, $\eD(Q)$ and $\dD(Q)$ (check the
identity of Proposition~\ref{prop:decomp}), and the three oracle bounds. Which one is the
smallest, and is this consistent with Proposition~\ref{prop:compare}?
\end{exercise}
\begin{exercise}
\textbf{(Application)} In Condorcet's model with $n=5$ independent voters of error $p=0.3$,
compute the exact risk of the vote, $\eD(Q)$, $\dD(Q)$ and the three oracle bounds. Check
the regime predicted by Proposition~\ref{prop:compare}.
\end{exercise}
\begin{exercise}
\textbf{(Application)} A posterior with $\KL(Q\|P)=1$ has an empirical Gibbs risk $0.15$ on
$m=1000$ labelled examples and an empirical disagreement $0.25$ on $m'=1000$ examples.
With $\delta=0.05$, compute $\overline r$, $\underline d$ and the PAC-Bayesian C-bound
\eqref{eq:cbound_pb}, the first-order certificate \eqref{eq:fo_pb} and the tandem
certificate \eqref{eq:tnd_pb} (the empirical joint error follows from
Proposition~\ref{prop:decomp}, assuming the disagreement was measured on the same
$1000$ labelled examples). Compare with the value of the C-bound computed at the empirical
values. What changes if $m'=10\,000$ unlabeled examples are available?
\end{exercise}
\begin{exercise}
\textbf{(Proof)} For voters with values in $\{-1,+1\}$ and a fixed example $(x,y)$, show that
$\P_{(h,h')\sim Q^2}\big(h(x)\neq h'(x)\big)=2W_Q(x,y)\big(1-W_Q(x,y)\big)$. Deduce that
$\E_\D M_Q^2=1-2\dD(Q)$ and $\RD(\GQ)=\eD(Q)+\frac12\dD(Q)$.
\end{exercise}
\begin{exercise}
\textbf{(Proof)} In Condorcet's model with $n=2k+1$ independent voters of error $p<\frac12$,
use Hoeffding's inequality (Lemma~\ref{lem:hoeffding}) to show that
$\RD(\BQ)\leq e^{-2n(1/2-p)^2}$. Evaluate this bound for $p=0.3$ and $n\in\{5,25,101\}$, and compare
it with the exact risks and with the C-bound.
\end{exercise}
\end{exercises}
